
%%%%%%%%%%%%%%%%%%%%%%%%%%%%%%%%%%%%%%%%%%%%%%%%%%%%%%%
%%%%%%%%%%%%%%%%%%%%%%%%%%%%%%%%%%%%%%%%%%%%%%%%%%%%%%%
%%%%%%%%%%%%%%%%%%%%%%%%%%%%%%%%%%%%%%%%%%%%%%%%%%%%%%%

\begin{frame}[t]{Introduction to covariate balance: What are we weighting for?\footnote{Pun borrowed from \textcite{solon:2015:weightingfor}}}
$$
\tarcol{\textrm{Target average response} =
\meantar \y_j} \approx \surcol{\meansur \w_i \y_i
= \textrm{Weighted survey average response }}
$$
We can't check this, because we don't observe $\tarcol{\y_j}$.  \pause But we can check whether:
$$
% \textrm{Target average regeressor} =
    \tarcol{\meantar \x_j} \eqcheck \surcol{\meansur \w_i \x_i}
% =    \textrm{Weighted survey average regressor}
$$

Weights that pass this check satisfy ``covariate balance'' for $\x$.

\pause

\vspace{1em}

You can check covariate balance for any weighting estimator,
and any function $\f(\x)$.

\textbf{Raking calibration weights} aim to exactly balance some set of regressors.


\pause

\begin{block}{Key question}
%
Covariate balance seems inherently linear. 

How can the there be an analogue of it for non-linear estimators?
%
\end{block}


\end{frame}




%%%%%%%%%%%%%%%%%%%%%%%%%%%%%%%%%%%%%%%%%%%%%%%%%%%%%%%
%%%%%%%%%%%%%%%%%%%%%%%%%%%%%%%%%%%%%%%%%%%%%%%%%%%%%%%
%%%%%%%%%%%%%%%%%%%%%%%%%%%%%%%%%%%%%%%%%%%%%%%%%%%%%%%

\begin{frame}[t]{Balance checks as local sensitivity}

\begin{block}{Key question}
%
Covariate balance seems inherently linear. 

How can the there be an analogue of it for non-linear estimators?
%
\end{block}

\onslide<2->{
\textbf{Key idea:}\\
The key idea is to convert the diagnostic into a \emph{local sensitivity analysis} of this form:

\begin{enumerate}
\item Assume your initial model was accurate
\item Select some perturbation your model should be able to capture
\item Use local sensitivity to detect whether the change is what you expect
\item If the change is not what you expect, either (1) or (2) was wrong
\end{enumerate}
}
%
\onslide<3->{
\vspace{2em}
One reason to balance $f(\x)$ is because we think
$\expect{}{\y \vert \x}$ might plausibly vary $\propto f(\x)$,
and want to check whether our estimator can capture this variability.

For example, linear regression balances regressors, and captures
variability in $\expect{}{\y \vert \x}$ that is proportional to the regressors.
}

\end{frame}

%%%%%%%%%%%%%%%%%%%%%%%%%%%%%%%%%%%%%%%%%%%%%%%%%%%%%%%
%%%%%%%%%%%%%%%%%%%%%%%%%%%%%%%%%%%%%%%%%%%%%%%%%%%%%%%
%%%%%%%%%%%%%%%%%%%%%%%%%%%%%%%%%%%%%%%%%%%%%%%%%%%%%%%

\begin{frame}[t]{Balance checks as local sensitivity}

One reason to balance $f(\x)$ is because we think
$\expect{}{\y \vert \x}$ might plausibly vary $\propto f(\x)$,
and want to check whether our estimator can capture this variability.

\only<2>{
\begin{block}{Covariate balance as local sensitivity (informal)}
    Pick a small $\delta > 0$ and an $\f(\cdot)$.  Define a \emph{new response variable} $\ytil$ such that
    $$
    \expect{}{\ytil \vert \x} = \expect{}{\y \vert \x} + \delta f(\x).
    $$
    We know the change this is supposed to induce in the target population.\\[1em]

    Covariate balance checks whether our estimators produce the same change.
    %
\end{block}
}
%
\only<3->{
\begin{block}{Covariate balance as local sensitivity (formal)}
    Pick a small $\delta > 0$ and an $\f(\cdot)$.  Define a \emph{new response variable} $\ytil$ such that
    $$
    \expect{}{\ytil \vert \x} = \expect{}{\y \vert \x} + \delta f(\x).
    $$
    We know the expected change this perturbation produces in the target distribution:
    $$
    \begin{aligned}
        \tarcol{
            \expect{}{\mu(\ytil) - \mu(\y) | \x} =
            \meantar \left(\expect{}{\ytil | \x_p} - \expect{}{\y | \x_p}\right) =
            \delta \meantar f(\x_j)}
    \end{aligned}
    $$
    Then, check whether your estimator $\muhat(\cdot)$ produces
    the same change for observed $\Ytil, \Ysur$:
    $$
    \begin{aligned}
        \underbrace{
            \tarcol{\muhat}(\Ytil) -
            \tarcol{\muhat}(\Ysur)
        }_{
            \substack{
                \text{Replace weighted averages} \\
                \text{with changes in an estimator}
            }
        }
        \overset{\textrm{\textbf{check}}}{\approx}
        \tarcol{
        \delta \meantar f(\x_j).}
    \end{aligned}
    $$
\end{block}
}


\end{frame}





%%%%%%%%%%%%%%%%%%%%%%%%%%%%%%%%%%%%%%%%%%%%%%%%%%%%%%%
%%%%%%%%%%%%%%%%%%%%%%%%%%%%%%%%%%%%%%%%%%%%%%%%%%%%%%%
%%%%%%%%%%%%%%%%%%%%%%%%%%%%%%%%%%%%%%%%%%%%%%%%%%%%%%%

\begin{frame}[t]{Balance checks as local sensitivity}


When $\tarcol{\muhat}(\cdot) = \muhatcw[\cdot]$,
we recover the standard covariate balance check.

$$
\begin{aligned}
    \underbrace{
        \muhatcw[\Ytil] - \muhatcw
    }_{
        \substack{
            \text{Replace weighted averages} \\
            \text{with changes in an estimator}
        }
    }
    ={}&
    \surcol{\meansur \w_i \ytil_i - \meansur \w_i \y_i}
    \\={}&
    \surcol{\meansur \w_i (\y_i + \f(\x_i)) - \meansur \w_i \y_i}
    \\={}&
    \surcol{\meansur \w_i \f(\x_i)}
    \\\eqcheck{}&
    \tarcol{
    \delta \meantar f(\x_j).}
\end{aligned}
$$


We will study $\tarcol{\muhat}(\cdot) = \muhatmrp[\cdot]$.

\end{frame}


%%%%%%%%%%%%%%%%%%%%%%%%%%%%%%%%%%%%%%%%%%%%%%%%%%%%%%%
%%%%%%%%%%%%%%%%%%%%%%%%%%%%%%%%%%%%%%%%%%%%%%%%%%%%%%%
%%%%%%%%%%%%%%%%%%%%%%%%%%%%%%%%%%%%%%%%%%%%%%%%%%%%%%%

\begin{frame}[t]{Balance checks for MrP}

Suppose I have
$\ytil$ such that $\expect{}{\ytil \vert \x} = \expect{}{\y \vert \x} + \delta f(\x)$.\\
Now I need to evaluate \surcol{$\muhatmrp[\Ytil] - \muhatmrp$}.
\pause

\textbf{Problem:} \surcol{$\muhatmrp[\cdot]$} is computed with MCMC.
%
\begin{itemize}
\item Each MCMC run typically takes hours, and
\item MCMC output is noisy, and \surcol{$\muhatmrp[\Ytil] - \muhatmrp$} may be small.
\end{itemize}
%
\pause
\textbf{Solution:} Use our local approximation, MrPlew!

\begin{block}{Balance informed sensitivity check with MrPlew:}
For a wide set of judiciously chosen $\f(\cdot)$, check
$$
\begin{aligned}
\muhatmrp[\Ytil] - \muhatmrp \approx{}&
    \surcol{\meansur \w_i^\mrp (\ytil_i - \y_i)}
    \\\approx{}&
    \underbrace{
        \delta  \surcol{\meansur \w_i^\mrp \f(\x_i)}
            \overset{\textrm{\textbf{check}}}{\approx}
        \delta \tarcol{\meantar f(\x_j).}
    }_{\textrm{What you actually check}}
\end{aligned}
$$
\end{block}

\end{frame}

%%%%%%%%%%%%%%%%%%%%%%%%%%%%%%%%%%%%%%%%%%%%%%%%%%%%%%%
%%%%%%%%%%%%%%%%%%%%%%%%%%%%%%%%%%%%%%%%%%%%%%%%%%%%%%%
%%%%%%%%%%%%%%%%%%%%%%%%%%%%%%%%%%%%%%%%%%%%%%%%%%%%%%%

\input{binary_y}

%%%%%%%%%%%%%%%%%%%%%%%%%%%%%%%%%%%%%%%%%%%%%%%%%%%%%%%
%%%%%%%%%%%%%%%%%%%%%%%%%%%%%%%%%%%%%%%%%%%%%%%%%%%%%%%
%%%%%%%%%%%%%%%%%%%%%%%%%%%%%%%%%%%%%%%%%%%%%%%%%%%%%%%

\begin{frame}[c]{Theory}

\textbf{When is the local approximation accurate?}

\begin{block}{Theorem: (sketch)}
    Take $\surcol{\ytil_i = \y_i + \delta \f(\x_i)}$.\\[1em]
    We state conditions for Bayesian hierarchical logistic regression under which
$$
\onslide<4->{\sup_{f \in \mathcal{F}}}
\abs{
    \muhatmrp[\Ytil] - \muhatmrp
    - \delta \surcol{\sumsur \w_i^\mrp \f(\x_i)}
}
\only<1>{ = \textrm{Small}}
\only<2->{ = O(\delta^2)}
\onslide<3->{\textrm{ as }N \rightarrow \infty}
$$
\onslide<4->{
    ...for a very broad class of $\mathcal{F}$.
    \footnote{$\mathcal{F}$ can be any Donsker class of measurable functions
    with uniformly bounded
    $\expect{}{\norm{\x} \f(\x)}_2^2 $.}
}
\end{block}

\onslide<4->{
    \textbf{Uniformity justifies searching for ``imbalanced'' $\f$.}
}

\onslide<5->{
The uniformity result builds on our earlier work on uniform
and finite--sample error bounds
for Bernstein--von Mises theorem--like results
\footcite{giordano:2024:bayesij,kasprzak:2025:laplace}.
}

\end{frame}


%%%%%%%%%%%%%%%%%%%%%%%%%%%%%%%%%%%%%%%%%%%%%%%%%%%%%%%%%%%%%%%%%
%%%%%%%%%%%%%%%%%%%%%%%%%%%%%%%%%%%%%%%%%%%%%%%%%%%%%%%%%%%%%%%%%
%%%%%%%%%%%%%%%%%%%%%%%%%%%%%%%%%%%%%%%%%%%%%%%%%%%%%%%%%%%%%%%%%



\begin{frame}{Covariate balance for primary effects}
\AlexanderImbalancePrimary{}
\end{frame}


\begin{frame}{Covariate balance for interaction effects}
\AlexanderImbalanceInteraction{}
\end{frame}




\begin{frame}[t]{Predictions}
    \AlexanderPredictionFigOne{}
\end{frame}


\begin{frame}[t]{Predictions and actual MCMC results}
    \AlexanderPredictionFigTwo{}
    \vspace{-3em}
    $$
    \begin{aligned}
        \textrm{Running ten MCMC refits:}\quad & \AlexRefitTimeHours \textrm{ hours} &
        \textrm{Computing approximate weights:}\quad &\AlexMrPawTimeSecs \textrm{ seconds}\\
    \end{aligned}
    $$
\end{frame}



%%%%%%%%%%%%%%%%%%%%%%%%%%%%%%%%%%%%%%%%%%%%%%%%%%%%%%%%%%%%%%%%%
%%%%%%%%%%%%%%%%%%%%%%%%%%%%%%%%%%%%%%%%%%%%%%%%%%%%%%%%%%%%%%%%%
%%%%%%%%%%%%%%%%%%%%%%%%%%%%%%%%%%%%%%%%%%%%%%%%%%%%%%%%%%%%%%%%%

% \begin{frame}[c]{Other uses}
% \centering


% Applying the idea to the observations negative weights measures
% \emph{non-monotonicity}.

% \wholeslidefig{
% \AlexanderWeightPlot{}
% }

% % \splitpagenoline{
% %     \AlexanderWeightPlot{}
% % }{
% %     \LaxWeightPlot{}
% % }
% \end{frame}
